\section{Additional Literature Review}
\label{app:additional_literature}

This appendix expands the related-work discussion of
Section~\ref{sec:related} along three axes that are most relevant to
score-debiased and other bias-corrected density estimators: sub-quadratic
approximate KDE algorithms, sample compression and coreset methods, and
KDE-as-primitive applications in modern machine learning.
We also discuss potential hybrids of these algorithmic directions with the
hardware-accelerated exact pipeline of Section~\ref{sec:method}.
The discussion is intentionally complementary to the experimental focus of
the main paper: the works below approximate either the kernel sum or the
training set, while our work makes the exact SD-KDE estimator practical at
scale.

\subsection{Sub-quadratic approximate KDE via locality-sensitive hashing}
\label{app:hbe_literature}

A line of work starting with the hashing-based estimator (HBE) framework of
\citet{charikar2017hbe} reduces the cost of evaluating a Gaussian kernel
sum from quadratic to sub-quadratic by sampling from locality-sensitive
hash buckets.
At a high level, HBE hashes both the training set and the query point into
LSH buckets whose collision probability is calibrated to the kernel, and
uses bucket cardinalities together with a small number of refinement
samples to estimate the kernel sum to a target relative error.
\citet{backurs2019space} improve the construction so that it requires only
linear preprocessing time and space while retaining the same query time,
and \citet{siminelakis2019rehashing} bring the scheme closer to practice
through adaptive sample-size selection and sharper data-dependent variance
bounds.
\citet{charikar2019multi} extend the construction beyond Gaussian-like
kernels to a much broader class of pairwise summations using a
multi-resolution hashing scheme.
A subsequent refinement \citep{charikar2020kde} casts the sub-quadratic KDE
problem as a density-constrained near-neighbor search and removes the
redundancy that arises in HBE when nearby points fall into different
buckets on opposite sides of the query.

The ACE method of \citet{luo2018ace}, which we evaluate against in
Section~\ref{app:ace}, is a practical instantiation of this family
specialized for streaming anomaly detection: rather than estimating the
kernel sum, it tracks per-bucket counts and uses them as proxies for local
density.
The closely related RACE sketch of \citet{coleman2020race} extends the
LSH-counting view to streaming KDE with explicit error guarantees.
ACE was a natural choice for our anomaly-detection comparison precisely
because it shares the high-level inverse-density-as-anomaly viewpoint of a
KDE-based scorer while bringing the sub-quadratic LSH cost.

These methods are complementary to ours rather than substitutes.
HBE-family algorithms approximate \emph{classical} KDE, so applying them
out of the box to SD-KDE would forfeit the bias-reduction guarantees that
motivate SD-KDE in the first place.
A natural hybrid direction (Section~\ref{app:future_work}) is to use a
sub-quadratic approximation only for the empirical-score pass and retain
exact streaming Tensor-Core evaluation for the final density.

\subsection{Sample compression, coresets, and kernel thinning}
\label{app:compression_literature}

A second line of work reduces the cost of kernel computations by replacing
the full training set with a much smaller weighted subset.
\citet{dwivedi2021kernel} introduce kernel thinning, which produces a
coreset of size $O(\sqrt{n})$ whose worst-case kernel-mean error matches
that of the full sample up to logarithmic factors, by halving the dataset
under a self-balancing process and selecting the half that better
preserves a target reproducing-kernel Hilbert space functional.
\citet{shetty2022compress} extend this construction with a meta-procedure
(Compress++) that reduces the total runtime to near-linear while inflating
the integration error by at most a constant factor, making kernel thinning
practical at substantially larger scales.

Compression-based methods affect a different cost axis from the one we
target.
They reduce $n$ before any kernel evaluations are performed, while we keep
$n$ fixed and accelerate the resulting $O(n^2)$ computation.
For SD-KDE specifically, kernel thinning could in principle be used to
reduce both the training set and the empirical-score sample, at the cost
of inflating bias and variance by amounts that have not yet been
characterized for score-debiased estimators.
We view a careful theoretical and empirical study of kernel thinning under
SD-KDE as out of scope for this paper but a promising follow-up.

\subsection{Higher-order bias correction in KDE}
\label{app:higher_order_literature}

The Laplace-corrected estimator of Section~\ref{sec:laplace_corrected}
belongs to the broader family of higher-order bias kernels.
Classical work in this area includes the higher-order kernel constructions
of \citet{fan1992bias} and the systematic comparison of
\citet{jones1997comparison}, both already discussed in the main related-work
section.
The novelty of our use of this construction is not in the kernel itself but
in identifying it as the right first-order surrogate for SD-KDE in the
limit where the empirical score is replaced by its leading-order Taylor
expansion, and in giving a fused GPU realization that retains the same
effective bias reduction at a fraction of the runtime cost.

\subsection{KDE as a primitive in modern ML systems}
\label{app:kde_attention_literature}

Although the present paper targets density estimation directly, KDE-style
sums also arise as computational primitives inside modern deep-learning
pipelines.
Self-attention with a softmax kernel can be written as a normalized
weighted sum that is closely related to a Gaussian KDE evaluated at the
query, and several recent works exploit this view to design
sub-quadratic attention approximations.
\citet{zandieh2023kdeformer} treat the attention denominator as a KDE and
use a sub-quadratic KDE algorithm in the spirit of the FOCS line above to
approximate it, obtaining provable spectral-norm guarantees and large
end-to-end speedups.
\citet{kacham2024polysketchformer} replace the softmax kernel with a
polynomial kernel and use polynomial-kernel sketches from randomized
numerical linear algebra to obtain a linear-time attention mechanism with
approximation guarantees.
Streaming-style LSH sketches in the spirit of RACE \citep{coleman2020race}
have similarly been adapted as the scoring primitive inside attention
modules.
FlashAttention~\citep{dao2022flashattention} is the closest analogue on the
hardware-aware side: it does not approximate the kernel sum but reorders
the exact computation to fit the GPU memory hierarchy and to expose
Tensor-Core GEMMs, which is the same high-level recipe we apply to SD-KDE.

These works underline that fast KDE-style primitives have downstream value
beyond density estimation.
The implementation pattern in this paper (matrix-multiplication
re-ordering, streaming accumulation, mixed Tensor-Core/FP32 work) is a
natural building block for any future architecture in which a
score-corrected or higher-order kernel sum is used in place of vanilla
attention or a vanilla KDE primitive.

\subsection{Hybrid algorithmic--hardware directions for future work}
\label{app:future_work}

The literature surveyed above suggests several hybrid directions that pair
algorithmic reductions with the hardware-accelerated exact pipeline of this
paper.

\paragraph{Approximate-score, exact-evaluation hybrids.}
The empirical score in SD-KDE is itself a sum of weighted kernels, and the
SD-KDE estimator is empirically robust to small score perturbations.
This makes it natural to plug a sub-quadratic approximate KDE algorithm,
such as those of
\citet{charikar2017hbe,backurs2019space,siminelakis2019rehashing,charikar2019multi,charikar2020kde},
into the score-evaluation step while retaining the exact Tensor-Core
kernel of Section~\ref{sec:method} for the final density evaluation.
Quantifying how approximate-score error propagates into the bias-reduction
guarantee of SD-KDE is an open theoretical question in this direction.

\paragraph{Spectral and grid-based score acceleration.}
For low-to-moderate dimensions, the score can also be computed by
evaluating the gradient of a smoothed density on a grid, with non-uniform
fast Fourier transforms (NUFFT) reducing the cost of the underlying sum.
Combining a NUFFT-based score with the exact Flash-SD-KDE evaluation kernel
is a promising hybrid for the dimensionalities we target, and is on our
roadmap.

\paragraph{Coreset-then-accelerate.}
A complementary hybrid uses kernel thinning
\citep{dwivedi2021kernel,shetty2022compress} or another coreset construction
to reduce $n$ before invoking the exact pipeline, trading a controlled
error increase for a square-root reduction in compute.
Because Flash-SD-KDE is already efficient at large $n$, the regime where
this trade is attractive likely lies at the very largest scales, where
even a hardware-accelerated exact $O(n^2)$ pass becomes prohibitive.

In all three directions, a fast and accurate exact SD-KDE primitive is
useful as the inner loop and as a ground-truth reference for evaluating
approximation error, which is the role this paper aims to fill.
